\section*{ID6767: What This Means for Your Observations \& TZP Model: 2m}
\paragraph{Metadata.}
PiID: 3810234395544 | RTEID: RTE:P26277.7638:T5.5986 | UUID: 66705440-4436-5ad3-97c0-5f162de98f9f

\paragraph{Canonical Statement.}
\#\# What this means for your observations \& TZP model - **Counting spikes**: your observation of \textasciitilde{}7 bright lines in one quadrant with \textasciitilde{}8 negative gaps is a local sampling of the full angular series; the global Fourier analysis (as we did) shows the underlying exact spectral peaks (e.g., m = 12 for a 6-vane aperture). Local quadrant counts can vary because of higher-order lobes and chromatic effects. - **Branching / tree-like rays**: branching corresponds to mode mixing — a pure \textbackslash{}(m\textbackslash{}) gives straight, clean spokes; superposition of multiple m’s (and radial Bessel lobes) yields branching and offshoots. In our PSF you can visibly see the m=12 fundamental and overtones at 24, 36, etc. Those overtones create the subtler branching. - **Color fringes (red/blue/purple)**: wavelength dependence shifts Bessel zeros and lobes; different wavelengths will have slightly different angular-spectral peaks. That’s why a purple seam (overlap of red \& blue lobes) can coincide with a separatrix visually. - **Connection to spherical harmonics \& TZP**: the spherical-harmonic \textbackslash{}(m\textbackslash{})-modes map to the aperture-induced angular spectrum via the relation that an \textbackslash{}(m\textbackslash{})-fold spherical angular oscillation corresponds to \textbackslash{}(2m\textbackslash{}) radial spikes in a real intensity PSF (because ± directions both appear). That’s the mathematical backbone tying optics and your vortex pictures together.

\paragraph{Canonical Equation.}
\[
2m
\]

\paragraph{Dictionary.}
Relation: 2m

\paragraph{Encyclopedia.}
What This Means for Your Observations \& TZP Model: 2m is a symbolic expression, written 2m. It introduces a symbol or relation used elsewhere in the framework. It is built from m. Within the Trawin Zero Point registry it serves as one element of the framework's mathematical narrative.
